\section{A sign-associativity criterion and a rank-nine exclusion}

All categories in this section are fusion categories over $\mathbb C$.
Let $R$ be a multiplicity-free fusion ring whose distinguished basis
$b_0=1,b_1,\ldots,b_{r-1}$ is self-dual. An \emph{admissible sign table}
is an associative integral multiplication table $T$ with unit $b_0$,
$T_{ij}^0=\delta_{ij}$, full symmetry of $T_{ij}^k$ in $i,j,k$, and
$|T_{ij}^k|=N_{ij}^k$. A basis-sign change has the form
$T_{ij}^k\mapsto s_i s_j s_k T_{ij}^k$, where $s_0=1$ and
$s_i\in\{1,-1\}$.

\begin{proposition}[Sign-associativity criterion]
If every admissible sign table of $R$ is a basis-sign change of $R$,
then every fusion categorification of $R$ admits a spherical structure.
Consequently, if no formal codegree $D$ of $R$ has $D/f$ an algebraic
integer for every formal codegree $f$, then $R$ has no categorification.
\end{proposition}
\begin{proof}
Let $\mathcal C$ be a hypothetical categorification and take its
canonical sphericalization $\widetilde{\mathcal C}$
\cite[Remark~3.1]{ENO}\cite[Section~5.1]{Ostrik3}.
The two lifts of each simple object differ by the sign object
$\varepsilon$, of spherical dimension $-1$. Every lift of a self-dual
simple is self-dual: its dual has the same nonzero spherical dimension,
whereas the other lift has the opposite dimension.
Choose a unit-preserving set of lifts. In
$K_0(\widetilde{\mathcal C})/(\varepsilon+1)$, their structure
constants form an admissible sign table. Self-duality makes this
quotient commutative, and Frobenius reciprocity gives full symmetry.
Multiplicity one makes each nonzero coefficient equal to $1$ or $-1$.

By hypothesis, tensoring some selected lifts with $\varepsilon$ makes
all nonzero quotient coefficients positive. The product of two such
lifts has exactly the selected lift in every nonzero channel: the
total multiplicity in the two lifts is one and their difference is
$+1$. Hence the selected lifts span a spherical fusion subcategory of
$\widetilde{\mathcal C}$. The forgetful functor restricts to a tensor
equivalence with $\mathcal C$, proving sphericality.

For a spherical categorification, its dimension $D$ is the formal
codegree of its dimension character. Every $D/f$ is an algebraic
integer by \cite[Theorem~2.13]{Ostrik3}. The last assertion follows.
\end{proof}

\begin{corollary}
The supplied ring $R9\_1\_41$, equivalently the AnyonWiki ring
$\mathrm{FR}^{9,0}_{38}$, has no fusion categorification over
$\mathbb C$.
\end{corollary}
\begin{proof}
Use the zero-based basis order in the supplied tensor. There are
$47$ nonzero unordered triples of non-unit labels. Basis-sign changes
allow the following eight coefficients to be set equal to $+1$:
\[
122,\quad 228,\quad 236,\quad 237,\quad245,\quad247,\quad256,\quad334.
\]
Here $ijk$ denotes $T_{ij}^{k}$. Independence of these eight gauge
conditions is an elementary rank-eight calculation over $\mathbb F_2$.
Thirty-six two-term consequences of associativity give the following
table; the remaining nonzero coefficients are all $+1$:
\[
\begin{array}{c|l}
a&335,347,366,446,455,567,777\\
b&336,344,357,456,477,555,667\\
c&337,346,355,444,457,566,677.
\end{array}
\]
The $b_4$ coefficient in $(b_3b_3)b_4=b_3(b_3b_4)$ then reads
\[
a+ab+c=3.
\]
Since $a,b,c\in\{1,-1\}$, this forces $a=b=c=1$.
Thus every admissible sign table is gauge-equivalent to $R$.
The supplementary certificate lists each of the $36$ source
associativity identities and checks every deduction by exact integer
arithmetic. No indicator normalization is used.

The characteristic polynomial of the regular Casimir matrix is
\[
(x-15)^6(x-2)(x^2-120x+1200).
\]
Thus the possible formal codegrees $D$ are $2$, $15$, and
$60\pm20\sqrt6$. The first two fail because $2/15$ and $15/2$
are nonintegral rationals. For the last two,
\[
N_{\mathbb Q(\sqrt6)/\mathbb Q}\!\left(\frac{60\pm20\sqrt6}{15}\right)
=\frac{16}{3},
\]
so the ratios are not algebraic integers. The proposition excludes
general categorification.
\end{proof}

\paragraph{Relation to the literature.}
Vercleyen--Slingerland, \emph{On Low-Rank Multiplicity-Free Complex
Fusion Categories}, Communications in Mathematical Physics
\textbf{407}, 231 (2026), Remark~4.3,
DOI \texttt{10.1007/s00220-026-05705-2}, record the pivotal exclusion
of this ring. The sign-associativity argument above supplies the
missing reduction from general categorification. It is a new
application of sphericalization to this particular tensor, not a
claim that the pivotality conjecture is known in general.

% Bibliography entries expected in the enclosing document:
% ENO: P. Etingof, D. Nikshych, V. Ostrik, On fusion categories,
% Ann. Math. 162 (2005), 581--642, arXiv:math/0203060.
% Ostrik3: V. Ostrik, Pivotal fusion categories of rank 3,
% Moscow Math. J. 15 (2015), 373--396, arXiv:1309.4822.
